\documentclass[a4paper,12pt]{article}

% 字体与页面
\usepackage[UTF8]{ctex}
\usepackage{fontspec}
\setmainfont{Times New Roman}

\usepackage{geometry}
\geometry{left=2.5cm, right=2.5cm, top=2.5cm, bottom=2.5cm}
\linespread{1.2} % n 倍行距
\setlength{\parskip}{0em} % 段间距紧凑

% 页眉页脚
\usepackage{fancyhdr}
\pagestyle{fancy}
\fancyhf{} % 清除页眉页脚
\cfoot{\arabic{page}} % 页脚序号
\renewcommand{\headrulewidth}{0pt} % 清除页眉页脚线
\setlength{\headheight}{15pt}

% 标题格式
\usepackage{titlesec}
\renewcommand{\thesection}{\chinese{section}}
\renewcommand{\thesubsection}{\arabic{section}.\arabic{subsection}}
\renewcommand{\thesubsubsection}{\arabic{section}.\arabic{subsection}.\arabic{subsubsection}}
\titleformat{\section}{\zihao{3}\heiti}{\thesection、}{0em}{}
\titleformat{\subsection}{\zihao{4}\heiti}{\thesubsection}{0.5em}{}
\titleformat{\subsubsection}{\zihao{4}\heiti}{\thesubsubsection}{0.5em}{}

% 代码环境
\usepackage{xcolor}
\usepackage{listings}
\usepackage{listingsutf8}
\usepackage{upquote}   % 直引号显示
\usepackage{subcaption}
\usepackage{booktabs}


\lstset{ 
  backgroundcolor=\color{gray!5},      % 浅灰背景
  frame=single,                        % 边框
  rulecolor=\color{gray!50},           % 边框颜色
  framerule=0.4pt,                     % 边框粗细
  framesep=6pt,                        % 边框与代码间距
  xleftmargin=0.5em, xrightmargin=0.5em, % 左右缩进
  aboveskip=0.8em, belowskip=0.8em,    % 上下间距
  basicstyle=\ttfamily\small,          % 等宽字体，小号
  numbers=left,                        % 显示行号
  numberstyle=\scriptsize\color{gray}, % 行号样式（更小、更淡）
  numbersep=8pt,                       % 行号和代码的间距
  breaklines=true,                     % 自动换行
  breakatwhitespace=false,             % 可在单词中间断行
  breakindent=0pt,                     % 换行后不缩进
  prebreak=\raisebox{0ex}[0ex][0ex]{\(\hookleftarrow\)\,}, % 换行前标记
  postbreak=\mbox{\space\(\hookrightarrow\)\space},        % 换行后标记
  tabsize=4,                           % Tab=4空格
  keepspaces=true,                     % 保留空格对齐
  showstringspaces=false,              % 不标记字符串里的空格
  keywordstyle=\bfseries\color{blue!70!black}, % 关键字 蓝+加粗
  commentstyle=\itshape\color{green!40!black}, % 注释 斜体墨绿
  stringstyle=\color{orange!90!black},         % 字符串 橙色
  literate=
    *{0}{{{\color{purple}0}}}{1}
     {1}{{{\color{purple}1}}}{1}
     {2}{{{\color{purple}2}}}{1}
     {3}{{{\color{purple}3}}}{1}
     {4}{{{\color{purple}4}}}{1}
     {5}{{{\color{purple}5}}}{1}
     {6}{{{\color{purple}6}}}{1}
     {7}{{{\color{purple}7}}}{1}
     {8}{{{\color{purple}8}}}{1}
     {9}{{{\color{purple}9}}}{1}
}

% 其他包
\usepackage[super, sort, compress]{cite}
\usepackage{graphicx, booktabs, tabularx}
\usepackage{amsmath}
\usepackage{amsfonts}
\usepackage{amssymb}
\usepackage{appendix}
\usepackage{caption}
\usepackage{longtable}
\usepackage{float}

\captionsetup{font=small}
\usepackage{enumitem}
\setlist[itemize]{itemsep=0.5ex, parsep=0ex, topsep=0.5ex}
\usepackage{array}
\newcolumntype{C}{>{\centering\arraybackslash}X}

% 信息
\title{}
\author{}
\date{}


%%%%%%%%%%%%%%% 正文 %%%%%%%%%%%%%%%
\begin{document}

\begin{center}
{\zihao{-2}\heiti 量化交易作业3}
\end{center}

\section{三种投资组合策略的分析}

\subsection{基本设定与数据处理}

\begin{itemize}
\item \textbf{数据：}工商银行、中国石油、万华化学、保利发展、贵州茅台 5 只股票的日频收盘价，1年期国债收益率（从财政部官网 https://gks.mof.gov.cn/gzsylqs/gzlssj/ 获取）；
\item \textbf{投资期限：}2010年1月1日至2017年12月4日，共7年；
\item \textbf{滚动调仓：}50天建模、月频调仓；
\item \textbf{数据预处理：}首先，读取数据，将0读取为缺失值np.nan，调整日期列的格式；
然后，由于万华化学在2017年12月4日到2018年6月4日之间停牌（清洗数据中发现），为了保证验证区间的完整，我截取了2017年12月4日之前的数据。
\end{itemize}

\subsection{平均收益率与协方差矩阵}

在每个月的最后一个交易日，用过去50个交易日的数据计算平均收益率和协方差矩阵，其中协方差矩阵的最小特征值序列如下图所示。
可以看出，最小特征值整体较小，多数位于$2\times10^{-5}$到$2\times10^{-4}$之间，整体差距不大，最大值也不到$3\times10^{-4}$。
且最小特征值均大于0，因此不需要担心协方差矩阵的正定性问题。

\begin{figure}[H]
\centering
\includegraphics[width=0.8\textwidth]{Pictures/Hw3_最小特征值时间序列.png}
\caption{最小特征值的序列}
\end{figure}


\subsection{三种投资组合策略的设定}

在每个月末调仓日，设 $w = (w_1, w_2, \dots, w_5)^T$ 为 5 只股票的权重，$\mu$ 为期望收益率，$\Sigma$ 为协方差矩阵，$\mathbf{1}$ 为全 1 的向量。
考虑A股市场的真实情况，加上禁止做空的约束，求解三种策略组合会变为优化问题：

\begin{itemize}
\item \textbf{全局最小方差组合}：只考虑风险资产。
\begin{equation*}
\begin{gathered}
\min_{w}\; w^T \Sigma w \\
\text{s.t.}\; w^T \mathbf{1} = 1,\quad 0 \le w_i \le 1
\end{gathered}
\end{equation*}

\item \textbf{均值-方差最优组合}：将目标收益率 $\mu_{target}$ 设定为 5 只股票期望收益率的平均值。
\begin{equation*}
\begin{gathered}
\min_{w}\; w^T \Sigma w \\
\text{s.t.}\; w^T \mu + \left(1-w^T \mathbf{1}\right) r_f = \mu_{target},\quad w^T \mathbf{1} \le 1,\quad 0 \le w_i \le 1
\end{gathered}
\end{equation*}

\item \textbf{二次效用最优组合}：设定投资者的风险厌恶系数 $\alpha = 4$。
\begin{equation*}
\begin{gathered}
\max_{w}\; w^T \mu + \left(1-w^T \mathbf{1}\right) r_f - \frac{\alpha}{2} w^T \Sigma w\\
\text{s.t.}\; w^T \mathbf{1} \le 1,\quad 0 \le w_i \le 1
\end{gathered}
\end{equation*}
\end{itemize}

之后可以用Python的scipy库求解上述优化问题，得到每个调仓日三种策略的权重。

\subsection{月度平均收益率、策略标准差}

求解月度平均收益率和标准差，如下表所示。
可以看出，最小方差组合并没有达到最小风险，反而达到了最高的收益率，和均值-方差最优组合相比，
体现初高风险高收益的特点；而二次效用最优组合在$\alpha=4$的条件下达到了最高的标准差，且收益率并没有比均值-方差组合高出很多，整体效果是最差的。

\begin{table}[H]
\centering
\caption{三种策略的月度指标}
\begin{tabular}{lccc}
\toprule
& 最小方差组合 & 均值-方差最优组合 & 二次效用最优组合 \\
\midrule
月度平均收益率 & 1.10\% & 0.85\% & 0.91\% \\
月度标准差     & 5.24\% & 3.32\% & 8.23\% \\
\bottomrule
\end{tabular}
\end{table}

\subsection{PnL曲线、年化收益率、标准差、夏普比率}

从PnL曲线可以看出，三种策略均实现了长期的正收益。
最小方差组合累计收益最高，中期波动较大，但在2016年后有最平稳的正向增长；
均值-方差组合虽然累计收益最低，但波动也是最小的，走势比较平缓；
二次效用最优组合波动最大，并且有连续多年为负收益，表现最差。


\begin{figure}[H]
\centering
\includegraphics[width=0.8\textwidth]{Pictures/Hw3_PnL曲线.png}
\caption{三种投资组合策略的累计 PnL 曲线}
\end{figure}

通过下表可以看出，年化收益率和标准差结果和月度的基本一致。
而夏普比显示，最小方差组合表现最好，二次效用最优组合表现最差，
也和我们对均值、标准差的分析结果相同。

\begin{table}[H]
\centering
\caption{年化收益率、标准差与夏普比率}
\begin{tabular}{lccc}
\toprule
& 最小方差组合 & 均值-方差最优组合 & 二次效用最优组合 \\
\midrule
年化收益率 & 14.73\% & 11.30\% & 12.12\% \\
年化标准差 & 19.10\% & 14.37\% & 30.94\% \\
夏普比率   & 0.6274  & 0.5956  & 0.3030  \\
\bottomrule
\end{tabular}
\end{table}

\subsection{权重分配变化和调仓}

如下图，最小方差组合的权重分配比较稳定，持有较多工商银行、中国石油和贵州茅台股票，每次调仓幅度较小，
而它们正是波动率最低的三只，体现了最小方差的目标。

\begin{figure}[H]
\centering
\includegraphics[width=0.8\textwidth]{Pictures/Hw3_min_var_资产权重.png}
\caption{最小方差组合资产权重变化}
\end{figure}

如下图，均值-方差组合的权重分配明显比最小方差组合变化更大，且多个时期持有无风险资产比例较大，
这可能是因为我们设定的目标收益率在变动，因此调仓会受到较大的影响，调仓幅度也比较大。

\begin{figure}[H]
\centering
\includegraphics[width=0.8\textwidth]{Pictures/Hw3_mv_资产权重.png}
\caption{均值-方差最优组合资产权重变化}
\end{figure}

如下图，二次最优组合的权重分配比较极端，长时间出现持有单一资产的情况，调仓幅度最大，
所以这也导致了波动率较大的结果。

\begin{figure}[H]
\centering
\includegraphics[width=0.8\textwidth]{Pictures/Hw3_utility_资产权重.png}
\caption{二次效用最优组合资产权重变化}
\end{figure}

\section{50天滚动的 $\beta$ 和 $R^2$ 分析}

\subsection{$\beta$ 系数和 $\beta$ 系数排序的变化}

每个交易日，用过去 50 个交易日的个股收益率和市场收益率进行滚动回归，得到每只股票 $\beta$ 系数的序列。
从下图 $\beta$ 的变化可以看出，保利发展（红线）和万华化学（绿线）的 $\beta$ 整体较高，多数时间在$>1$的水平；
其余三者的 $\beta$ 整体较低，多数时间在$<1$的水平。

\begin{figure}[H]
\centering
\includegraphics[width=0.8\textwidth]{Pictures/Hw3_beta_series.png}
\caption{$\beta$ 系数的时间序列}
\end{figure}

进一步，下表的描述性统计结果也可以验证上述结论：保利发展和万华化学的 $\beta$ 系数均值和分位点都是最高的，并且显著高于其他股票；
另外，五只股票 $\beta$ 的标准差大小接近，说明波动率差别不大。

\begin{table}[H]
\centering
\caption{$\beta$ 系数的描述性统计}
\begin{tabular}{lccccc}
\toprule
& 工商银行 & 中国石油 & 万华化学 & 保利发展 & 贵州茅台 \\
\midrule
均值 & 0.4574 & 0.5758 & 1.1437 & 1.2454 & 0.6447 \\
标准差  & 0.2346 & 0.2469 & 0.2652 & 0.2560 & 0.2823 \\
最小值  & -0.3257 & -0.0738 & 0.0833 & -0.4906 & -0.1433 \\
25\% & 0.3275 & 0.4249 & 0.9534 & 1.0869 & 0.4599 \\
50\% & 0.4443 & 0.5309 & 1.1440 & 1.2460 & 0.6255 \\
75\% & 0.5514 & 0.6718 & 1.2825 & 1.4399 & 0.7938 \\
最大值  & 1.3028 & 1.4806 & 2.1571 & 1.7591 & 1.6488 \\
\bottomrule
\end{tabular}
\end{table}

从下图的 $\beta$ 的排序变化来看，万华化学和保利发展多数时间排前两名，这里和前文的结论是一致的。

\begin{figure}[H]
\centering
\includegraphics[width=0.8\textwidth]{Pictures/Hw3_beta_rank.png}
\caption{$\beta$ 系数排名的序列}
\end{figure}

\subsection{当期 $\beta$ 系数和前一期的 $\beta$ 系数的 $R^2$}

下面对两种 $R^2$ 进行计算和对比：第一种是用当期数据得到的 $\beta$ 计算当期的 $R^2$；
第二种是用前一期数据得到的 $\beta$ 计算当期的 $R^2$。前者时样本内，后者是样本外，通过对比，我们可以分析出 $\beta$ 是否具有预测能力。
下图以贵州茅台股票为例，可以看出两种 $R^2$ 在绝大多数时间是重合的，说明 $\beta$ 系数稳定性很好，因此具备预测能力。

\begin{figure}[H]
\centering
\includegraphics[width=0.8\textwidth]{Pictures/Hw3_R2.png}
\caption{样本内拟合 $R^2$ 与样本外预测 $R^2$ 的时间序列对比}
\end{figure}

如下表，从数值上来看，两种 $R^2$ 的平均值也非常接近，进一步验证了 $\beta$ 系数具有稳定性和预测能力。

\begin{table}[H]
\centering
\caption{两种 $R^2$ 的平均值对比}
\begin{tabular}{lccccc}
\toprule
& 工商银行 & 中国石油 & 万华化学 & 保利发展 & 贵州茅台 \\
\midrule
样本内$R^2$ & 0.3457 & 0.4553 & 0.4412 & 0.4622 & 0.2275 \\
样本外$R^2$  & 0.3440 & 0.4540 & 0.4400 & 0.4610 & 0.2260 \\
\bottomrule
\end{tabular}
\end{table}

\end{document}